\titledquestion{File Systems}

\textbf{File Systems}

Consider the structure of a Linux file system, which is composed of \textit{directories}
and \textit{files}.

A directory is a structure which may contain further directories and files. A file
is just a standalone piece of data.

We may represent it via the following SML type:\footnote{Which is (almost) isomorphic to
a \code{string rose}, for those of you keeping track at home.}
\begin{codeblock}
  datatype filesys =
      File of string
    | Dir of string * filesys list

  type filesys = (string * data) list
  and data = File | Dir of filesys
\end{codeblock}

A \textit{file system} is a value of type \code{filesys} such that for each
\code{Dir (d, ds)}, no \code{File} or \code{Dir} in \code{ds} shares the
same name.

So for instance, a file system with the following structure:

\begin{minipage}{\linewidth} % Use a minipage to control the width
  \dirtree{%
  .1 root.
  .2 A.
  .3 foo.sml.
  .3 bar.txt.
  .2 B.
  .3 baz.sml.
  .2 C.
  }
\end{minipage}

would be represented by the following SML value:
\begin{codeblock}
  Dir ("root",
    [ Dir ("A", [ File "foo.sml", File "bar.txt" ]),
      Dir ("B", [ File "baz.sml"]),
      Dir ("C", [])
    ]
  )
\end{codeblock}

\newpage

You are given an implementation of the \code{cd} function, which allows
moving through the file system, according to the following specification:

\spec
  {cd}
  {\code{filesys -> string list -> filesys option}}
  {\code{true}}
  {\code{cd fs path} takes the moves indicated by the path \code{path},
  and returns \code{SOME fs'} if that ends up at a file system,
  and \code{NONE} if it doesn't.}

So for the above filesystem, if it was to be bound to \code{fs}, we would
have that:
\begin{itemize}
  \item \code{cd fs []} $\eeq$ \code{SOME fs}
  \item \code{cd fs ["B"]} $\eeq$ \code{SOME (Dir ("B", [ File "baz.sml" ]))}
  \item \code{cd fs ["root"]} $\eeq$ \code{NONE}
\end{itemize}

The given implementation is:

\begin{codeblock}
  fun cd fs [] = fs
    | cd (File _) (s::rest) = NONE
    | cd (Dir (dir, ds)) (s::rest) =
        List.foldl
          (fn (_, SOME res) => SOME res
            | (File _, _) => NONE
            | (Dir (dir', ds'), _) =>
                if s = dir' then
                  cd (Dir (dir', ds')) rest
                else
                  NONE
          )
          NONE
          ds
\end{codeblock}

\begin{parts}
  \part[6]\hfill

    Assume that any directory in a given file system has $k$
    many entries in it, and it is roughly balanced.

    Write and solve a recurrence for the \textbf{work} of \code{cd} on
    such a file system, in terms of the \textbf{depth} of the file system.
    The depth of a file system is, intuitively, the maximum number of
    nested directories you can go into.

    \begin{solutionorbox}[18em]
      $W_{\code{cd}}(0) = c_0$
      $W_{\code{cd}}(d) = O(k) + W_{\code{cd}}(d - 1)$

      So ultimately, $O(kd)$.
    \end{solutionorbox}

    \begin{EnvUplevel}
      Consider the following function on file systems:

      \begin{codeblock}
        fun countdirs (File _) = 0
          | countdirs (Dir (_, dirs)) =
              List.map countfiles dirs
              |> List.foldl (op+) 1
      \end{codeblock}
    \end{EnvUplevel}

  \part[5]\hfill

    Assume that any directory in a given file system has $k$
    many entries in it, and it is roughly balanced.

    Write but \textbf{do not solve} a recurrence for the \textbf{work} of
    \code{countdirs} on such a file system, in terms of the \textbf{nodes} of
    the file system. The nodes of a file system is, intuitively, the total
    number of files and directories in it.

    \begin{solutionorbox}[18em]
      $W_{\code{cd}}(0) = c_0$
      $W_{\code{cd}}(n) = O(k) + k \cdot W_{\code{cd}}(n / k)$
    \end{solutionorbox}

  \newpage

  \part[6]\hfill

    Assume that any directory in a given file system has $k$
    many entries in it, and it is roughly balanced.

    Write and solve a recurrence for the \textbf{span} of \code{countdirs} on
    such a file system, in terms of the \textbf{nodes} of the file system.
    The nodes of a file system is, intuitively, the total number of
    files and directories in it.

    \textbf{For this problem, and only this problem, pretend that \code{List.map}
    runs each call in parallel.}

    \begin{solutionorbox}[20em]
      $S_{\code{cd}}(0) = c_0$
      $S_{\code{cd}}(n) = O(k) + S_{\code{cd}}(n / k)$

      So ultimately, $O(k \cdot \log_k n)$.
    \end{solutionorbox}





  % \part[5]\hfill

  %   Assume that any directory in a given file system has, at maximum, $k$
  %   many entries in it.

  %   Write a recurrence for the \textbf{span} of \code{cd} on
  %   such a file system, in terms of the \textbf{depth} of the file system.
  %   The depth of the file system is, intuitively, the maximum number of
  %   nested directories you can go into.

  %   You may find useful the fact that $1 + k + k^2 + ... + k^n$ is on the order
  %   of $O(k^n)$, for non-zero $k$.

  %   \begin{solutionorbox}[8em]
  %     $S_{\code{cd}}(0) = c_0$
  %     $S_{\code{cd}}(d) = O(k) + \cdot S_{\code{cd}}(d - 1)$
  %   \end{solutionorbox}


  \newpage

  \part[10]\hfill

    Implement \code{absolutize}, which replaces all of the names in the file
    system with its absolute path. (which should start with a "/"). For
    instance, running \code{absolutize} on the above example file system should
    give you
    \begin{codeblock}
      Dir ("/root",
        [ Dir ("/root/A", [ File "/root/A/foo.sml",
                            File "/root/A/bar.txt" ]),
          Dir ("/root/B", [ File "/root/B/baz.sml"]),
          Dir ("/root/C", [])
        ]
      )
    \end{codeblock}

    \spec
      {absolutize}
      {\code{filesys -> filesys}}
      {\code{true}}
      {\code{absolutize fs} replaces each file and directory's name with
      the absolute path required to get there}

    \begin{solutionorbox}[30em]
      \begin{codeblock}
        fun mkPath strs =
          List.foldl (fn (x, acc) =>
            acc ^ "/" ^ x
          ) "" strs

        fun absolutize' cur (File s) = File (mkPath (s :: cur))
          | absolutize' cur (Dir (dir, ds)) =
              let
                val new_cur = dir :: cur
              in
                Dir (mkPath new_cur, List.map (absolutize new_cur) ds)
              end

        fun absolutize = absolutize' []
      \end{codeblock}
    \end{solutionorbox}
\end{parts}

